\section{The Geometry of the Substrate}
\label{sec:geometry}

The substrate is a regular tessellation of a negatively curved space, grown rather than placed.
This section treats it in full, because it is what the rest of the model rests on.

\subsection{Coxeter groups and Schl\"afli symbols}

A regular tessellation is the orbit of a single tile under a group generated by reflections in the
tile's faces. The mathematics of such reflection groups is the theory of Coxeter groups \citep{coxeter1973regular}. A Coxeter group is specified by the angles between its mirrors, and for a regular tessellation
those angles are encoded in a Schl\"afli symbol, a short list of integers. The symbol $\{p\}$ is a
$p$-gon. The symbol $\{p,q\}$ is the tiling in which $q$ regular $p$-gons meet at every vertex. The
symbol $\{p,q,r\}$ is the honeycomb in which $r$ cells of type $\{p,q\}$ meet at every edge, and so
on in higher dimensions. The only data are whole numbers, and the curvature of the space the
tessellation lives in is fixed by them.

\paragraph{The curvature condition.}
For the tiling $\{p,q\}$, the quantity $1/p + 1/q$ decides the geometry. If it exceeds $1/2$ the
tiling closes up on a sphere, a Platonic solid. If it equals $1/2$ the tiling is the flat plane,
giving exactly the three Euclidean cases $\{3,6\}$, $\{4,4\}$, $\{6,3\}$. If it is less than $1/2$
the tiling lives in the hyperbolic plane, and there are infinitely many such, one for every pair
$(p,q)$ with $1/p + 1/q < 1/2$. The smallest hyperbolic case at $q = 3$ is $\{7,3\}$, since $\{6,3\}$
is already flat. So seven is the first integer that curves space at three tiles to a vertex, which
is why the heptagrid is the natural minimal illustration of the whole framework.

\subsection{Regular tilings of the hyperbolic plane}

\begin{figure}[h]
\centering
\includegraphics[width=0.5\textwidth]{vibe-mesh-73.pdf}
\caption{The heptagrid $\{7,3\}$, the regular tiling of the hyperbolic plane by seven-sided cells,
three to a vertex, drawn in the Poincar\'e disc. Each tile is a vibe and each shared edge a note.
Tiles look smaller toward the rim only because the disc compresses the infinite hyperbolic plane
into a finite picture. In the metric of the space every tile is the same size, and the rim is
infinitely far away, so there is always more room at the edge.}
\label{fig:heptagrid}
\end{figure}

Figure~\ref{fig:heptagrid} shows the heptagrid in the Poincar\'e disc, the conformal model that
fits the whole infinite hyperbolic plane inside a circle. Three features matter for the model.
First, the number of tiles within graph distance $n$ of a centre grows exponentially in $n$, the
defining signature of negative curvature, so the crystal has unbounded room to grow. Second, the
tiling is exactly regular, every tile congruent, yet it has no global preferred direction, because
in hyperbolic space there is no way to carry a direction from one point to a distant one without
ambiguity. Third, the disc picture is a faithful map of the relations, not a container: the physics
is in which tiles touch, and any conformal redrawing leaves that untouched. The testbed builds
these tilings by Margenstern's splitting method \citep{margenstern2003theory}, with each node assigned a
fixed number of children by its type, reproducing the exact ball-growth law of the tiling.

\subsection{Regular honeycombs of hyperbolic space, and the dodecagrid}

\begin{figure}[h]
\centering
\includegraphics[width=\textwidth]{vibe-mesh-534.png}
\caption{The dodecagrid $\{5,3,4\}$, the order-4 dodecahedral honeycomb of hyperbolic
three-space, the model's actual substrate. Dodecahedral cells fill negatively curved 3D space,
four meeting at every edge. This is the three-dimensional analogue of the heptagrid, harder to draw
but the same kind of object. Time is the growth of this crystal.}
\label{fig:dodecagrid}
\end{figure}

In three dimensions a regular tessellation is a honeycomb $\{p,q,r\}$ filling a curved space with
polyhedral cells. The model's actual substrate is the one shown in Figure~\ref{fig:dodecagrid}, the
order-4 dodecahedral honeycomb $\{5,3,4\}$, called the dodecagrid. Its cells are dodecahedra, twelve
pentagonal faces each, and four cells meet at every edge, a packing that only closes up in
hyperbolic three-space. The dodecagrid is to three dimensions what the heptagrid is to two: an
exact, regular, negatively curved crystal with exponential volume growth and no global preferred
direction. The testbed constructs it as a graph with a Poincar\'e-ball embedding
(\texttt{hyperbolicDodecagrid}), and uses it wherever the genuinely three-dimensional substrate is
needed, for navigation and for the mass-hierarchy geometry below.

\subsection{The classification that fixes the dimension}

The spatial dimension is not assumed in the model. It is forced, twice over.

\paragraph{The dimension window (P62).}
Regular honeycombs with finite cells, which is what a substrate of finite, locally connected vibes
must be, exist as compact regular hyperbolic honeycombs only in a narrow band of dimensions. In the
hyperbolic plane there are infinitely many. In hyperbolic three-space there are exactly four (the
two orders of the dodecahedral honeycomb among them). In hyperbolic four-space there are exactly
five. In five dimensions and above there are none: hyperbolic space can no longer be filled with
finite regular cells at all. The testbed reproduces this known classification exactly (counts of
$\infty, 4, 5, 0, 0, \ldots$ for dimensions $2,3,4,5,6,\ldots$), so the substrate dimension is a
computed constraint, not a free parameter.

 The universe must be two-, three-, or four-dimensional in
space simply to be built of finite cells.

\paragraph{Three among the window (P68, P69).}
Within that window, three is selected. P68 shows that only three spatial dimensions support stable,
closed gravitational orbits, the inverse-square law that lets atoms and solar systems hold together:
two dimensions precesses and four is unstable, the classical Ehrenfest argument \citep{ehrenfest1917way}
recovered on the discrete substrate. P69 measures the dimension that emerges from pure grown
connectivity, with no embedding put in, and finds it converges to a sharp value as the mesh grows.
So the dimension is constrained to a window of three values and then selected to one, by stability
and by the growth itself.

\subsection{Why curvature buys Lorentz invariance}

Discreteness is widely expected to break Lorentz invariance, because a regular flat lattice has
preferred axes: its photon speed becomes direction- and energy-dependent, which gamma-ray-burst
timing tightly forbids. The substrate escapes this, and the reason is geometric, not statistical.
In a flat lattice a direction can be carried unchanged across the whole lattice, so the axes are
globally defined and physical. In a curved space a direction cannot be carried unambiguously from
one point to a distant one, so the link directions of a regular hyperbolic crystal look, to any
local observer, isotropic, like a random foam with no preferred frame. P27 measures this directly:
a lattice has an order-one anisotropy that grows with energy, while the hyperbolic substrate's
anisotropy is consistent with zero and falls with size.

 The lesson is that curvature, not disorder,
is what buys Lorentz safety, so the substrate can be an exact ordered crystal and still look smooth
and relativistic from within. P82 turns this into numbers: the substrate predicts no first-order
Lorentz violation and passes the Fermi-LAT bound that excludes a lattice.

\subsection{Addressing and navigation}

Because the crystal is a tree-like hyperbolic graph, it can be addressed and navigated locally. P42
gives the two-dimensional case, following Margenstern's theory of these tilings \citep{margenstern2003theory}:
each cell carries a Zeckendorf address, a sum of non-consecutive Fibonacci numbers written as a
binary string with no two adjacent ones, and a signal routes between any two cells by local address
comparison. P76 gives the three-dimensional case on the dodecagrid:
each cell's address is its position in the hyperbolic embedding, and greedy routing, stepping each
time to the neighbour whose address is closest to the target, delivers every pair tested at the
shortest path. Greedy routing works precisely because the geometry is hyperbolic, the same curvature
that buys Lorentz safety.

\subsection{Deterministic eternal growth}

Finally, the crystal is not placed but grown. P83 builds a deterministic automaton that adds cells
one at a time at the frontier, breadth-first, by the fixed splitting rule, with no randomness and no
rebuilding. The growth is append-only, so once a cell is added its links never change, the interior
frozen as the past and the activity all at the leading edge that is the present. It is resumable, so
growing in many small steps gives exactly the same crystal as one large step, which lets it continue
forever. And it is faithful, matching the statically built tiling cell for cell. Crucially the
geometry emerges from the growth: the ball-growth ratio converges on its own to the pentagrid's
golden-ratio law $(3+\sqrt 5)/2 \approx 2.618$, so the negative curvature is a consequence of the
growth rule rather than an input. This is the eternal-growth substrate the model needs: a crystal
that is also a process.
